\section{Companion App and Health}
\label{sec:companion-health}

\subsection{Companion App}
\label{sec:companion}

The system ships with a first-party companion app for Android and iOS that
bridges phone and desktop into a single coherent experience. It is built with
Tauri 2.0 using the same Rust backend, the same design token system, and the
same UI component library as the desktop. The companion app looks and feels
like the OS because it is built from the same parts.

The protocol is custom. KDE Connect was evaluated and rejected: its protocol
design would obstruct the deeper integration features planned for later tiers,
and building on top of it would mean working around its assumptions rather than
designing for the actual requirements.

\paragraph{Feature tiers.}
Features are organized into three trust tiers established at pairing time.
Table~\ref{tab:companion-tiers} shows the breakdown. The distinction between
\texttt{paired} and \texttt{trusted} exists because Tier 2 features involve
access to graph context and health data, which require a higher trust
level than file transfer and notification mirroring. Tier 2 is restricted
to the official app; a third-party client that implements the pairing protocol
only gets Tier 1.

\begin{table}[htbp]
\caption{Companion app feature tiers}
\label{tab:companion-tiers}
\begin{tabularx}{\linewidth}{llX}
\toprule
\textbf{Tier} & \textbf{Trust level} & \textbf{Features} \\
\midrule
1 & \texttt{paired}  &
  File Drop; clipboard sync (manual trigger only, not automatic);
  notification mirroring; media playback control \\
2 & \texttt{trusted} &
  Handoff (continue reading, browsing, or editing across devices);
  Unified Timeline (cross-device activity, local only);
  graph presence (phone sends active app and content context);
  health data forwarding \\
3 & Planned          &
  Wearable data pipeline; extended context sharing \\
\bottomrule
\end{tabularx}
\end{table}

\paragraph{Protocol and transport.}
Discovery uses mDNS service advertisement on a shared WLAN. Without a shared
network, the desktop and phone locate each other via BLE presence broadcast
carrying the device ID. Transport is selected by availability at connection
time, shown in Table~\ref{tab:companion-transport}.

\begin{table}[htbp]
\caption{Companion app transport selection}
\label{tab:companion-transport}
\begin{tabularx}{\linewidth}{lX}
\toprule
\textbf{Transport} & \textbf{Use case} \\
\midrule
HTTPS over LAN  & Primary transport; full feature set \\
WiFi Direct     & Large transfers without a shared router \\
Bluetooth       & Small data only: clipboard, notifications, control commands \\
\bottomrule
\end{tabularx}
\end{table}

Pairing uses a Trust on First Use model. Both devices generate an Ed25519
keypair on first launch. The first connection uses TLS with self-signed
certificates plus a six-digit confirmation code displayed on the desktop and
confirmed on the phone. After pairing, reconnection is automatic with no
further confirmation. Unpair deletes the key on both sides with no trace.

\paragraph{Auto-connect.}
Paired devices broadcast BLE presence in the background. When a known device
is detected, the system connects automatically using the best available
transport. No user action is required after initial pairing. Android supports
full background BLE advertising reliably. iOS imposes CoreBluetooth Background
Mode restrictions: auto-connect works but is less reliable and carries a
marginally higher battery cost on the iPhone side.

\paragraph{Theme synchronization.}
Because the companion app uses the same design token system as the desktop,
the active theme propagates to the phone automatically: accent color, dark and
light mode, typography. This is a structural side effect of the shared token
infrastructure, not a feature that was separately implemented.

\paragraph{iOS limitations.}
iOS imposes constraints that cannot be worked around. There is no persistent
background BLE scanning, which makes auto-connect less reliable and requires
more manual reconnects than on Android. WiFi Direct is not available to
third-party iOS apps. USB-C is inaccessible to third-party apps without
MFi certification. Health data forwarding via HealthKit is feasible with
explicit user permission.

Tier 1 works well on iOS. Tier 2 is functional but degraded. iOS is explicitly
a secondary platform and this limitation is communicated honestly in the app
and in documentation rather than obscured.

\subsection{Health and Wellbeing}
\label{sec:health}

The system includes a first-party health data component: a local daemon that
aggregates, stores, and optionally analyzes health and wellbeing data from
connected sources. This is not a fitness application. It is health data
infrastructure: an open, local equivalent of Apple Health or Google Fit, built
as a system component with the same privacy guarantees as the rest of the
platform.

\paragraph{Data architecture.}
Health data lives in a dedicated local store managed by a \emph{Health Daemon},
separate from the Graph Daemon. Data never leaves the device without an
explicit user action. Clinical data uses FHIR R4 (HL7 Fast Healthcare
Interoperability Resources) where applicable, enabling import from hospital and
clinic systems, which are increasingly FHIR-capable across the EU, and export
to FHIR-compatible tools without conversion. Non-clinical data such as steps,
sleep stages, heart rate variability, and nutrition uses an open JSON schema
with documented field definitions. There is no proprietary format; full export
is always available.

\paragraph{Data sources.}
The initial release supports manual input and phone sensors forwarded via the
Companion App: steps, GPS, and screen time. The plugin-based architecture is
designed from the start to accommodate further sources without structural
changes: wearables via the Companion App, hospital and clinic FHIR import, and
third-party applications via a documented Health API gated by the standard
permission system.

\paragraph{Wearable pipeline.}
Wearables connect to the system through the Companion App rather than directly,
using the wearable manufacturer's existing phone SDK on the phone side. The
data path is:

\begin{lstlisting}[language=bash, caption={Wearable data path}]
Wearable -> Companion App (phone, manufacturer SDK)
         -> OS Health Daemon (via Companion protocol)
\end{lstlisting}

No direct wearable protocol implementation is required in the OS. The Health
Daemon's plugin interface is designed to accept forwarded wearable data from
the start, so that adding wearable sources later does not require structural
changes. This is planned architecture, not initial release scope.

\paragraph{Graph integration.}
Health data integration with the Knowledge Graph is strictly opt-in and split
into three explicit levels, all of which default to off.
Table~\ref{tab:health-graph} lists the levels. Each level requires the
previous one; the Settings UI enforces this and does not allow incompatible
states. The AI layer's health-adjacent features are gated entirely on these
toggles. With all three off, the health module is a pure local viewer and
aggregator with no AI involvement whatsoever.

\begin{table}[htbp]
\caption{Health graph integration levels}
\label{tab:health-graph}
\begin{tabularx}{\linewidth}{lX}
\toprule
\textbf{Level} & \textbf{What it enables} \\
\midrule
1 - Health Nodes     &
  Health events appear as standalone nodes in the Knowledge Graph \\
2 - Correlations     &
  Graph links health data to activity data by time window
  (e.g.\ sleep quality correlated with focus session length) \\
3 - Pattern Recognition &
  Full cross-domain analysis across health, productivity, and
  communication data \\
\bottomrule
\end{tabularx}
\end{table}

\paragraph{Settings and permissions.}
Health has its own dedicated section in Settings rather than being buried under
general privacy settings. Controls include per-source enable and disable for
each data source independently; per-data-type enable and disable so the user
can, for example, share step data with an app but not sleep data; the graph
integration toggles from Table~\ref{tab:health-graph} presented as a grouped
sub-section with clear explanations of what each level enables; a configurable
data retention period per data type; full FHIR and JSON export; and full
deletion with a confirmation step.

Third-party applications request health data access per data type through the
standard permission system, but health-specific confirmation dialogs are more
explicit about data sensitivity than generic permission prompts. A fitness app
requesting sleep data gets a dialog that names the data type and explains
clearly that this is health information, not a generic file access request.
